\documentclass[12pt,a4paper]{article}
\usepackage[margin=1in]{geometry}
\usepackage{amsmath}
\usepackage{amssymb}
\usepackage{array}
\usepackage{booktabs}
\usepackage{xcolor}
\usepackage{fancyhdr}
\usepackage{graphicx}
\usepackage{setspace}

\onehalfspacing
\pagestyle{fancy}
\fancyhf{}
\fancyhead[L]{CPE 415 -- Digital Image Processing}
\fancyhead[R]{Mid Term SP23 -- Solutions}
\fancyfoot[C]{\thepage}
\renewcommand{\headrulewidth}{0.5pt}

\title{\textbf{Mid Term Examination -- Spring 2023} \\ \textbf{Complete Solutions}}
\author{CPE 415: Digital Image Processing \\ Instructor: Dr. Ikramullah Khosa}
\date{}

\begin{document}

\maketitle

\section*{Solution to Question 1}

\subsection*{Part A: Linear Interpolation}

Linear interpolation computes values between existing data points using the formula:
\begin{equation*}
f(x_0 + \Delta x) = f(x_0) + \Delta x \cdot [f(x_0 + 1) - f(x_0)]
\end{equation*}

For interpolation at $\frac{1}{4}$ pixel, $\Delta x = \frac{1}{4}$.

Between pixels 3 and 10 (at position $0.25$):
\begin{align*}
f(0.25) &= 3 + \frac{1}{4}(10 - 3) \\
&= 3 + 1.75 = 4.75 \\
&\approx 5 \text{ (rounded)}
\end{align*}

Between pixels 10 and 2 (at position $1.25$):
\begin{align*}
f(1.25) &= 10 + \frac{1}{4}(2 - 10) \\
&= 10 - 2 = 8 \\
&\approx 8
\end{align*}

\textbf{Result:} The interpolated values at quarter-pixel intervals are 5 and 8.

\subsection*{Part B: CCD Chip Resolution}

The resolving power of a camera depends on the spatial frequency capabilities determined by the imaging geometry.

Given parameters:
\begin{itemize}
\item CCD chip dimensions: $14 \times 14$ mm
\item CCD resolution: $2048 \times 2048$ pixels
\item Pixel pitch: $\frac{14 \text{ mm}}{2048} = 0.00684$ mm/pixel
\item Object distance: 0.5 m
\item Focal length: 35 mm
\end{itemize}

The magnification is:
\begin{equation*}
M = \frac{f}{s_o - f} = \frac{35 \text{ mm}}{500 \text{ mm} - 35 \text{ mm}} \approx 0.0757
\end{equation*}

The effective spatial sampling at the object plane:
\begin{equation*}
\text{Pixel size at object} = \frac{0.00684 \text{ mm}}{0.0757} \approx 0.0904 \text{ mm}
\end{equation*}

The maximum resolvable frequency (line pairs per mm):
\begin{equation*}
\text{Resolution} = \frac{1}{2 \times 0.0904} \approx 5.53 \text{ line pairs/mm}
\end{equation*}

\textbf{Answer:} The camera can resolve approximately 5.5 line pairs per mm.

\vspace{1.0cm}
\section*{Solution to Question 2}

\subsection*{Part A: First and Second Derivatives}

The first and second derivatives are computed using finite difference approximations:

First derivative (forward difference):
\begin{equation*}
\frac{\partial f}{\partial x}\bigg|_{x_i} \approx f(x_{i+1}) - f(x_i)
\end{equation*}

Second derivative (central difference):
\begin{equation*}
\frac{\partial^2 f}{\partial x^2}\bigg|_{x_i} \approx f(x_{i+1}) - 2f(x_i) + f(x_{i-1})
\end{equation*}

Given sequence: 4, 5, 6, 7, 7, 0, 6, 5, 1, 1

First derivative:
\begin{equation*}
\text{First} = [1, 1, 1, 0, -7, 6, -1, -4, 0]
\end{equation*}

Second derivative:
\begin{equation*}
\text{Second} = [0, 0, -1, -7, 13, -7, -3, 4]
\end{equation*}

\subsection*{Part B: Histogram Equalization}

The histogram equalization transformation is:
\begin{equation*}
s_k = \text{round}\left[(L-1) \sum_{j=0}^{k} p_r(r_j)\right]
\end{equation*}

where $L = 8$ (for 3-bit images).

Computing the cumulative probabilities and transformations:

\begin{center}
\begin{tabular}{|c|c|c|c|c|c|}
\hline
$r_k$ & $p_r(r_k)$ & $\sum p_r$ & $7 \times \sum p_r$ & $s_k$ & Mapping \\
\hline
0 & 0.19 & 0.19 & 1.33 & 1 & $0 \to 1$ \\
1 & 0.25 & 0.44 & 3.08 & 3 & $1 \to 3$ \\
2 & 0.21 & 0.65 & 4.55 & 5 & $2 \to 5$ \\
3 & 0.16 & 0.81 & 5.67 & 6 & $3 \to 6$ \\
4 & 0.08 & 0.89 & 6.23 & 6 & $4 \to 6$ \\
5 & 0.06 & 0.95 & 6.65 & 7 & $5 \to 7$ \\
6 & 0.03 & 0.98 & 6.86 & 7 & $6 \to 7$ \\
7 & 0.02 & 1.00 & 7.00 & 7 & $7 \to 7$ \\
\hline
\end{tabular}
\end{center}

The equalized histogram is:

\begin{center}
\begin{tabular}{|c|c|c|}
\hline
\textbf{Equalized Level} & \textbf{Source Levels} & \textbf{Probability} \\
\hline
0 & -- & 0.00 \\
1 & 0 & 0.19 \\
2 & -- & 0.00 \\
3 & 1 & 0.25 \\
4 & -- & 0.00 \\
5 & 2 & 0.21 \\
6 & 3, 4 & 0.24 \\
7 & 5, 6, 7 & 0.11 \\
\hline
\end{tabular}
\end{center}

\subsection*{Part C: Laplacian and Unsharp Masking Equivalence}

The Laplacian operator is defined as:
\begin{equation*}
\nabla^2 f = \frac{\partial^2 f}{\partial x^2} + \frac{\partial^2 f}{\partial y^2}
\end{equation*}

For a slowly varying smooth image, the Laplacian represents the difference between a pixel and its local average. Consequently:
\begin{equation*}
f(x, y) - \nabla^2 f(x, y) \approx f(x, y) - [f(x, y) - \bar{f}(x, y)] = \bar{f}(x, y)
\end{equation*}

Therefore:
\begin{equation*}
f(x, y) - \nabla^2 f(x, y) \sim f(x, y) - \bar{f}(x, y)
\end{equation*}

This equivalence demonstrates that subtracting the Laplacian performs the same enhancement operation as unsharp masking, which subtracts a blurred version from the original image.

\subsection*{Part D: Nyquist Sampling Rate for Checkerboard}

For a checkerboard pattern with square dimensions of $0.5 \times 0.5$ mm, each square represents one complete cycle of the spatial pattern.

The fundamental frequency of the checkerboard pattern:
\begin{equation*}
f_{\max} = \frac{1}{0.5 \text{ mm}} = 2 \text{ cycles/mm}
\end{equation*}

By the Nyquist sampling theorem, the minimum sampling rate must be at least twice the maximum frequency:
\begin{equation*}
f_s \ge 2 \times f_{\max} = 2 \times 2 = 4 \text{ samples/mm}
\end{equation*}

\textbf{Answer:} The minimum sampling rate required is 4 samples per mm.

\vspace{2.0cm}
\noindent\rule{\textwidth}{0.5pt}
\begin{center}
\textit{End of Solutions}
\end{center}

\end{document}
